\section{Lecture 3 (September 8th)}
\begin{defi}
(Grassmann numbers) Grassman numbers are numbers, contrarily to complex numbers, that describe fermions. Grassmann variables are tuplets 
\[\{\theta_1,\ldots ,\theta_{n}\}\]
that satisfy
\[\{\theta _{i},\theta _{j}\}=\theta _{i}\theta _{j}+\theta _{j}\theta _{i}=0\]
The Grassmann algebra is defined as
\[\mathrm{\Lambda }(v)={\bm C}\oplus V\oplus V^2\oplus \ldots \oplus V^{n}\]
where $V$ is an $n$-dimensional vector space with the basis $\{\theta _{i}\}$ over the field of complex numbers. Basic operations are defined as
\begin{itemize}
	\item[(i)] Addition $\theta _{i}+\theta _{j}$
	\item[(ii)] Complex multiplication $c\theta _{i}$ ($c\in {\bm C}$)
	\item[(iii)] Product $\theta_{i}\theta _{j}$
\end{itemize}
The general element of the Grassmann algebra is called the Grassmann number,
\[\begin{align*}
	f(\theta )=&f_{0}+\sum ^{n}_{i_{1}=1}f_{i_1}\theta _{i_1}+\dfrac{1}{2}\sum _{i_1,i_2=1}^{n}f_{i_1i_2}\theta _{i_1}\theta _{i_2}+\ldots \\
	=& \sum ^{n}_{k=0}\dfrac{1}{k!}\sum ^{n}_{i_1,\ldots ,i_{k}=1}f_{i_1\ldots i_{k}}\theta _{i_1}\ldots \theta _{i_{k}}
\end{align*}\]	
\end{defi}
\vspace{2ex}
\begin{defi}
(Differentiation of Grassmann variables) 
	\[\dfrac{\partial \theta _{j}}{\partial \theta _{i}}=\delta _{ij}\qquad \dfrac{\partial {\bf 1}}{\partial \theta _{i}}=0\qquad \dfrac{\partial }{\partial \theta _{i}}(\theta _{j}\theta _{k})=\delta _{ij}\theta _{k}   \]
\end{defi}
\vspace{2ex}
\begin{defi}
(Integration of Grassmann variables) 
	\[\int d\theta \,f(\theta )=\dfrac{\partial }{\partial \theta }f(\theta ) \]
Notice that once we change the integration measure from $\theta _{i}'=A_{ij}\theta _{j}$,
\[\begin{align*}
	\int d\theta _{1}\ldots d\theta _{n}\,f(\theta )=&\dfrac{\partial }{\partial \theta_1}\ldots \dfrac{\partial }{\partial \theta _{n}}f(\theta )  \\
	=&A_{i_{1}1}\ldots A_{i_{n}n}\dfrac{\partial }{\partial \theta _{i_1}'}\ldots \dfrac{\partial }{\partial \theta _{i_{n}}'}f(A^{-1}\theta ') \\
	=&\varepsilon _{i_{1}\ldots i_{n}}A_{i_11}\ldots A_{i_{n}n}\dfrac{\partial }{\partial \theta _{1}'}\ldots \dfrac{\partial }{\partial \theta _{n}'}f(A^{-1}\theta ')\\
	=&\mathop{\mathrm{det}}(A)\int d\theta _{1}'\ldots d\theta _{n}'\,f(A^{-1}\theta ')
\end{align*}\]
which gives us
\[d\theta_1'\ldots d\theta _{n}'=(\mathrm{\mathop{det}}A)^{-1}d\theta _{1}\ldots d\theta _{n}\]
as
\[\dfrac{\partial }{\partial \theta _{i}}=\dfrac{\partial \theta _{j}'}{\partial \theta _{i}}\dfrac{\partial }{\partial \theta _{j}'}   \]
\end{defi}
\vspace{2ex}
\begin{defi}
(Gaussian integral wth respect to two sets of Grassman variables)
\[\begin{align*}
	&\int d\bar{\theta }_{1}d\theta _{1}\ldots d\bar{\theta }_{n}d\theta _{n}\,e^{-\sum _{i,j=1}^{n}\bar{\theta }_{i}M_{ij}\theta _{j}}\\
	&=\int d\bar{\theta}_{1}d\theta _{1}'\ldots d\bar{\theta }_{n}d\theta _{n}'(\mathop{\mathrm{det}}M)e^{o\sum ^{n}_{i=1}\bar{\theta }_{i}\theta _{i}'}\\
	&=\mathop{\mathrm{det}}M\prod_{i=1}^{n}\int d\bar{\theta }_{i}d\theta _{i}'\,e^{-\bar{\theta }_{i}\theta _{i}'}\\
	&=\mathop{\mathrm{det}}M\prod^{n}_{i=1}\int d\bar{\theta }_{i}d\theta _{i}'(1-\bar{\theta}_{i}\theta _{i}')=\mathop{\mathrm{det}}M
\end{align*}\]
(where $M_{ij}\theta _{j}=\theta '_{i}$)
\end{defi}
\vspace{2ex}
\begin{defi}
(Functional derivatives) Let $X$ be a space of admissible functions $y:[a,b]\rightarrow {\bm R}$ like the Banach space $L^2([a,b])$ or $C^{1}([a,b])$. A functional is a mapping $F:X\rightarrow {\bm R}$ defined as
\[\dfrac{\delta F[f]}{\delta f(x_0)}=\lim _{\varepsilon \rightarrow 0}\dfrac{F[f(x)+\varepsilon \delta(x-x_0) ]-F[f(x)]}{\varepsilon }\]
\end{defi}
\vspace{2ex}
\begin{ex}
(Examples of functional derivatives) We have
\begin{itemize}
\item[(i)] $\int f^{n}$ gives
\[\dfrac{\delta F[f]}{\delta f(z)}=n(f(z))^{n-1}\]
\item[(ii)] $\int (f')^{n}$ gives
\[\dfrac{\delta F[f]}{\delta f(z)}=-nf^{(n)}(z)\]
\item[(iii)] $\int dx\, K(y,x)f(x)$ gives
\[\dfrac{\delta F[f]}{\delta f(z)}=K(y,z)\]
\item[(iv)] $f(y)$ gives
\[\dfrac{\delta F[f]}{\delta f(z)}=\delta (y-z)\]
\end{itemize}
\end{ex}
\vspace{2ex}
\section{Classical Field Theory}
\begin{thm}
(Lagrangian formalism) Classically, the action $S$ of a system is written as
\[S[x,\dot{x}]=\int ^{t_{f}}_{t_{i}}dt\,L(x_{i},\dot{x}_{i})\]
For a field, we write
\[S[\phi ,\dot{\phi }]=\int ^{t_{f}}_{t_{i}}dt\,L[\phi ,\dot{\phi }]\]
where $\phi =\phi (t,\vec{x})=\phi (x)$. The variation of this action is given as
\[\begin{align*}
	\delta S[\phi ,\dot{\phi }]=\int dt\,\int  d^3\vec{x}\,\Big[\dfrac{\delta L}{\delta \phi (x)}\delta \phi (x)+\dfrac{\delta L}{\delta \dot{\phi }(x)}\delta \dot{\phi }(x)\Big]
\end{align*}\]
Through the application of Hamilton's principle (the variation of the action is zero for an arbitrary variation of the field such that $\delta \phi (t_{i},\vec{x})=\delta \phi (t_{f},\vec{x})=0$) we obtain
\[\begin{align*}
	0=&\int d^{4}x\,\Big(\dfrac{\delta L}{\delta \phi (x)}\delta \phi (x)+\dfrac{\delta L}{\delta \dot{\phi }(x)}\dfrac{\partial }{\partial t}\delta \phi (x) \Big)\\
	=&\int d^{4}x\,\Big(\dfrac{\delta L}{\delta \phi (x)}-\dfrac{\partial }{\partial t}\dfrac{\delta L}{\delta \dot{\phi }(x)} \Big)\delta \phi (x)+\int d^{4}x\,\dfrac{\partial }{\partial t}\Big(\dfrac{\delta L}{\delta \dot{\phi }(x)}\delta \phi (x)\Big) 
\end{align*}\]
where the second boundary terms vanish to zero. The Lagrangian in "local" relativistic classical field theory leads to the concept of a Lagrangian density.
\[S[\phi ,\dot{\phi }]=\int dt\,L[\phi ,\dot{\phi }]=\int d^{4}x\,\mathcal{L}(\phi (x),\dot{\phi }(x),\nabla \phi (x))=\int d^{4}x\,\mathcal{L}(\phi (x),\partial _{\mu }\phi (x))\]
Alternative expressions in local classical field theory include
	\[\dfrac{\delta L}{\delta \phi }=\dfrac{\partial \mathcal{L}}{\partial \phi }-\nabla \Big(\dfrac{\partial \mathcal{L}}{\partial (\nabla \phi )} \Big)\qquad \dfrac{\delta L}{\delta \dot{\phi }}=\dfrac{\partial \mathcal{L}}{\partial \dot{\phi }}  \]
and the Euler-Lagrangian equations as
\[0=\dfrac{\partial \mathcal{L}}{\partial \phi }-\partial _{\mu }\Big(\dfrac{\partial \mathcal{L}}{\partial (\partial _{\mu }\phi )} \Big) \]
as
\[\begin{align*}
	0=&\int d^{4}x\,\Big[\dfrac{\partial \mathcal{L}}{\partial \phi }\,d\phi +\dfrac{\partial \mathcal{L}}{\partial (\partial _{\mu }\phi )}\partial _{\mu }(\delta \phi ) \Big]\\
	=&\int d^{4}x\,\Big(\dfrac{\partial \mathcal{L}}{\partial \phi }-\partial _{\mu }\Big(\dfrac{\partial \mathcal{L}}{\partial  (\partial_{\mu }\phi  )} \Big)\delta \phi +\int d^{4}x\,\partial _{\mu }\Big(\dfrac{\partial \mathcal{L}}{\partial (\partial _{\mu }\phi )}\delta \phi \Big) \Big)
\end{align*}\]
assuming that $\phi ,\partial _{\mu }\phi \rightarrow 0	$ fast enough as $|\vec{x}|\rightarrow \infty $. 
\end{thm}
\vspace{2ex}
\begin{thm}
(Hamiltonian formulation) The canonically conjugate field is defined as
\[\pi(x)=\dfrac{\delta L}{\delta \dot{\phi }(x)}=\dfrac{\partial \mathcal{L}}{\partial \dot{\phi }(x)} \]
The Hamiltonian is then defined as
	\[H[\phi ,\pi ]=\int d^3\vec{x}\,\Big(\pi (x)\dot{\phi }(x)-\mathcal{L}\Big)=\int d^3x\,\mathcal{H}\]
The Eular-Lagrange equations then become Hamilton's equations of motion
	\[\dot{\phi }(x)=\dfrac{\delta H}{\delta \pi (x)}\qquad \dot{\pi }(x)=-\dfrac{\delta H}{\delta \phi (x)}\]
\end{thm}
\vspace{2ex}

